% Visual illustration of Proposition 1. Requires \usepackage{tikz}.
\begin{tikzpicture}[x=1cm,y=1cm,
  vtx/.style={circle,inner sep=0pt,minimum size=4.2mm,font=\scriptsize\bfseries,text=white},
  lab/.style={font=\small},
  hdr/.style={font=\small\bfseries},
  bar/.style={line width=2.4pt,line cap=round}]
\definecolor{cA}{HTML}{2A78D6}  % residues 1-4
\definecolor{cB}{HTML}{EB6834}  % residues 5-8
\definecolor{fillg}{HTML}{D9D9D9}

% ---------- helpers ----------
% \sqr{x}{y}{color}{first label}{filled? 0/1}
\newcommand{\sqr}[5]{%
  \begin{scope}[shift={(#1,#2)}]
    \ifnum#5=1
      \fill[fillg] (0,0)--(0.9,0)--(0.9,0.9)--(0,0.9)--cycle;
      \draw[gray!70,thick] (0,0)--(0.9,0.9) (0.9,0)--(0,0.9);
    \fi
    \draw[black!70,thick] (0,0)--(0.9,0)--(0.9,0.9)--(0,0.9)--cycle;
    \pgfmathtruncatemacro{\a}{#4}\pgfmathtruncatemacro{\b}{#4+1}
    \pgfmathtruncatemacro{\c}{#4+2}\pgfmathtruncatemacro{\d}{#4+3}
    \node[vtx,fill=#3] at (0,0) {\a};   \node[vtx,fill=#3] at (0.9,0) {\b};
    \node[vtx,fill=#3] at (0.9,0.9) {\c}; \node[vtx,fill=#3] at (0,0.9) {\d};
  \end{scope}}
% \clq{x}{y}{color}{first label}: small filled 4-clique (the compact cluster C)
\newcommand{\clq}[4]{%
  \begin{scope}[shift={(#1,#2)}]
    \fill[fillg] (0,0)--(0.7,0)--(0.7,0.7)--(0,0.7)--cycle;
    \draw[black!70,thick] (0,0)--(0.7,0)--(0.7,0.7)--(0,0.7)--cycle (0,0)--(0.7,0.7) (0.7,0)--(0,0.7);
    \pgfmathtruncatemacro{\a}{#4}\pgfmathtruncatemacro{\b}{#4+1}
    \pgfmathtruncatemacro{\c}{#4+2}\pgfmathtruncatemacro{\d}{#4+3}
    \node[vtx,fill=#3] at (0,0) {\a};   \node[vtx,fill=#3] at (0.7,0) {\b};
    \node[vtx,fill=#3] at (0.7,0.7) {\c}; \node[vtx,fill=#3] at (0,0.7) {\d};
  \end{scope}}

% ---------- column headers ----------
\node[hdr] at (1.1,4.35) {$K_0$};
\node[hdr] at (4.75,4.35) {$K_0\cup K_1$};
\node[hdr] at (8.3,4.35) {$K_1$};
\node[hdr] at (11.3,4.35) {framewise $\mathrm{Dgm}_1$};
\node[hdr] at (14.6,4.35) {zigzag barcode};

% ================= Trajectory X (same residues keep the loop) =================
\node[lab,anchor=east] at (-0.4,3.05) {$X$};
\sqr{0}{2.6}{cA}{1}{0}    \clq{1.5}{2.7}{cB}{5}
\sqr{3.6}{2.6}{cA}{1}{0}  \clq{5.1}{2.7}{cB}{5}
\sqr{7.2}{2.6}{cA}{1}{0}  \clq{8.7}{2.7}{cB}{5}
\draw[->,thick,gray] (2.5,3.05)--(3.3,3.05);
\draw[<-,thick,gray] (6.1,3.05)--(6.9,3.05);
\node[font=\scriptsize,gray!60!black] at (4.75,2.15) {square $1\!-\!4$ is still a loop};

% ================= Trajectory X' (labels exchanged at t=1) =================
\node[lab,anchor=east] at (-0.4,0.75) {$X'$};
\sqr{0}{0.3}{cA}{1}{0}    \clq{1.5}{0.4}{cB}{5}
\sqr{3.6}{0.3}{cA}{1}{1}  \sqr{5.2}{0.3}{cB}{5}{1}
\sqr{7.2}{0.3}{cB}{5}{0}  \clq{8.7}{0.4}{cA}{1}
\draw[->,thick,gray] (2.5,0.75)--(3.3,0.75);
\draw[<-,thick,gray] (6.4,0.75)--(6.95,0.75);
\node[font=\scriptsize,gray!60!black] at (4.75,-0.15) {both squares filled: $H_1=0$};

% ---------- framewise diagram (identical for X and X') ----------
\begin{scope}[shift={(10.3,0.3)}]
  \draw[->] (0,0)--(2.2,0);
  \node[font=\scriptsize,anchor=north east] at (2.2,-0.05) {birth};
  \draw[->] (0,0)--(0,3.4);
  \node[font=\scriptsize,anchor=south] at (0,3.4) {death};
  \draw[gray,dashed] (0,0)--(2.1,3.2);
  \fill[black] (0.8,1.9) circle (2.4pt);
  \node[font=\scriptsize,right] at (0.85,1.9) {$(a,\,a\sqrt2)$};
  \node[font=\scriptsize,align=center] at (1.1,-0.75) {identical full diagrams\\at every frame};
  \node[font=\tiny,align=center] at (1.1,0.6) {square feature shown\\other points omitted};
\end{scope}

% ---------- zigzag barcodes ----------
\begin{scope}[shift={(13.6,0)}]
  \foreach \x/\t in {0/{$K_0$},1/{$\cup$},2/{$K_1$}} \node[font=\scriptsize] at (\x,-0.45) {\t};
  \foreach \x in {0,1,2} {\draw[gray!40] (\x,-0.25)--(\x,3.7);}
  \draw[bar,cA] (0,3.05)--(2,3.05);
  \node[font=\scriptsize] at (1,2.6) {$X$: one class persists};
  \fill[cA] (0,0.95) circle (3.2pt);
  \fill[cB] (2,0.55) circle (3.2pt);
  \node[font=\scriptsize] at (1,1.45) {$X'$: death, then birth};
\end{scope}
\end{tikzpicture}
